% FORMALIZACION_NUEVA / CERTIFICADO_NUEVO, 10-09-2026.
% Composition of the effective operations already assembled in 30e and 31.
% Does not modify route selection, boundary, couplings, or
% action normalization; does not evaluate a cosmological amplitude s_0.
% Owners and scope: desarrollo/COMPOSICION_METRICA_REV06.md.
% Independent check: pruebas/verificar_eliminacion_torsional_no_lineal.py.

\section{Torsional elimination and metric response of the minimal extension}
\label{vii:sec:eliminacion-torsional-no-lineal}

The minimal extension provides a field, a response weight, and a
current through a single variational operation. Its transported
incidence depends on the connection through the route products in
\eqref{vii:eq:derivada-incidencia-rutas}. That dependence is preserved when
the current is composed with the Cartan--Holst equation. The connection
is therefore eliminated from the complete action of that
extension. The affine case of
\ref{vii:thm:respuesta-cuadratica} is contained as an exact
specialization of the following construction.

\subsection{Variational composition at a cellular level}
\label{vii:subsec:composicion-torsional-celular}

Fix a finite level, its routes, the transport representation, and
the cellular pairing of
\ref{vii:thm:componentes-minimo-conexion}. Let \(p\) be real
coordinates of the cotetrad and material fields in a declared local
trivialization, and let \(k\in\mathbb R^d\) be the coordinates of
contorsion \(K\). The collection \(p\mapsto e(p)\) preserves
nondegeneracy of the cotetrad. The connection acting on each route is
\begin{equation}
 \omega(p,k)=\mathring\omega(e(p))+K(k),\qquad
 F(p,k)=S_{\min}[e(p),\omega(p,k),\Psi(p)]
       =\mathfrak a(p,k)E(p,k).
 \label{vii:eq:accion-minima-en-contorsion}
\end{equation}
The spaces and bases in this expression are transported to fixed
spaces before differentiation. We work on an open set where the incidence
\(\mathsf C(p,k)\) is surjective and the collections are \(C^2\).
This additional regularity permits examination of the action Hessian.
The couplings \(G_{\rm int},c_{\rm int},\gamma_{\rm HMT}\)
preserve the chart and variation convention of
\ref{vii:subsec:dominio-corriente-cartan}.

Let \(Q(p,k)\) be the integral, or the declared cellular realization, of
the quadratic form
\eqref{vii:eq:forma-cuadratica-contorsion}. In real bases,
\begin{equation}
 Q(p,k)=\tfrac12 k^{\mathsf T}\mathsf A(p)k,
 \qquad \mathsf A(p)^{\mathsf T}=\mathsf A(p).
 \label{vii:eq:cuadratica-celular-contorsion}
\end{equation}
The operator \(\mathsf A\) comes from the gravitational action and
its pairing, whereas the Gramian
\(\mathsf C\mathsf C^*\) comes from the extension problem.
Positivity of the latter guarantees existence of the minimizing
field; the gravitational form preserves its own signature.
Boundary terms are maintained according to
\eqref{vii:eq:borde-contorsion}; the following interior formulas
apply with admissible variations annihilating their contribution, or
with the corresponding boundary completion.

Define the response covector by
\begin{equation}
 d_kF(p,k)[h]= -\tfrac12 h^{\mathsf T}s(p,k),
 \qquad s(p,k)=\mathsf M_{\rm cel}(p,k)\sigma(p,k).
 \label{vii:eq:covector-torsional-minimo}
\end{equation}
Here \(s\) and \(\sigma\) preserve the distinct types in
\eqref{vii:eq:componentes-minimo-conexion}: the former contains the
coefficients of the differential, and the latter the coordinates of the
current in the material pairing. When the basis in \(k\)
changes, that pairing is transported as well. Differentiation of
\(s=\mathsf M_{\rm cel}\sigma\) includes derivatives of both
factors when the pairing depends on the parameter.

\begin{proposition}[Composed equation of extension and contorsion]
\label{vii:prop:ecuacion-torsional-no-lineal}
Interior stationarity of the action with respect to contorsion is
\begin{equation}
 \mathcal R(p,k):=\mathsf A(p)k-\tfrac12s(p,k)=0.
 \label{vii:eq:ecuacion-torsional-no-lineal}
\end{equation}
Each component of \(s\) is calculated from the same extension:
\begin{equation}
 \begin{split}
 s_j={}&4\mathfrak a\chi\operatorname{Re}
       \langle y,(\partial_{k_j}\mathsf C)f_{\min}
                         -\partial_{k_j}b\rangle\\
 &-2\mathfrak a(\partial_{k_j}\chi)\langle b,y\rangle
       -2(\partial_{k_j}\mathfrak a)E.
 \end{split}
 \label{vii:eq:seleccion-variacional-corriente-k}
\end{equation}
When the cellular pairing realizes the connection equation of
\ref{vii:thm:ecuacion-cartan-holst}, this equation is equivalent to composing
\(\sigma(p,k)\) with the Holst inverse, the inverse of
\(\mathsf L_e\), and reconstruction of contorsion:
\begin{equation}
 K=\mathsf C_e^{\rm tor}\,\mathsf L_e^{-1}
           \bigl(\kappa c\,\mathsf P_\gamma^{-1}
             \sigma[e,\mathring\omega(e)+K,\Psi]\bigr).
 \label{vii:eq:composicion-implicita-cartan}
\end{equation}
The symbol \(\mathsf C_e^{\rm tor}\) denotes exclusively the
map \(T\mapsto K\) of
\eqref{vii:eq:contorsion-componentes}, with fixed cotetrad.
\end{proposition}

\begin{proof}
The differential of \(Q+F\) in the arbitrary direction \(h\) is
\[
 d_k(Q+F)[h]
   =h^{\mathsf T}\bigl(\mathsf A k-\tfrac12s\bigr).
\]
Its vanishing in every direction proves
\eqref{vii:eq:ecuacion-torsional-no-lineal}.
Theorem~\ref{vii:thm:componentes-minimo-conexion}, applied to
\(\partial_{k_j}\), gives
\eqref{vii:eq:seleccion-variacional-corriente-k}. With \(p\) fixed,
\(\delta\omega=\delta K\). The Cartan equation is transformed
successively through \eqref{vii:eq:inversa-holst},
\eqref{vii:eq:torsion-explicita}, and
\eqref{vii:eq:contorsion-componentes}; each of these operations
is invertible on its domain. This proves equivalence with
\eqref{vii:eq:composicion-implicita-cartan} when both expressions
use the same material and cellular realization.
\end{proof}

The composition preserves the connection on both sides. Routes
traversing several cells may couple the current components
between them. Pointwise reconstruction of \(T\) from a given
current and joint solution of
\eqref{vii:eq:ecuacion-torsional-no-lineal} are, in that case,
successive operations. This finite formulation keeps explicit
the route representation entering the functional.

\subsection{Local existence of a stationary branch}
\label{vii:subsec:rama-torsional-local}

\begin{theorem}[Local continuation of the torsional solution]
\label{vii:thm:rama-torsional-local}
Suppose \((p_0,k_0)\) belongs to the preceding open set,
\(\mathcal R(p_0,k_0)=0\), and the total Hessian
\begin{equation}
 \mathsf H_0
 =\mathsf A(p_0)+\partial_k^2F(p_0,k_0)
 =\mathsf A(p_0)-\tfrac12\partial_ks(p_0,k_0)
 \label{vii:eq:hessiano-total-torsional}
\end{equation}
is invertible. There exist neighbourhoods of \(p_0\) and \(k_0\) in which
a unique \(C^1\) branch, \(k=k_*(p)\), satisfies
\eqref{vii:eq:ecuacion-torsional-no-lineal} and passes through \(k_0\).
On that branch,
\begin{equation}
 \partial_{p_i}k_*
 =-\bigl[\mathsf A+\partial_k^2F\bigr]^{-1}
     \left[(\partial_{p_i}\mathsf A)k_*
                      +\partial_{p_i}\partial_kF\right].
 \label{vii:eq:derivada-rama-torsional}
\end{equation}
\end{theorem}

\begin{proof}
Gramian inversion is \(C^2\) on the full-rank open set.
Thus \(F\) is \(C^2\), and \(\mathcal R\) is
\(C^1\). Its derivative with respect to \(k\), at \((p_0,k_0)\), is
\(\mathsf H_0\). The implicit function theorem provides
the branch, its regularity, and uniqueness in a neighbourhood of the point.
Existence can also be seen through the map
\[
 \mathcal T_p(k)=k-\mathsf H_0^{-1}\mathcal R(p,k).
\]
Its derivative with respect to \(k\) vanishes at \((p_0,k_0)\).
In a sufficiently small ball its norm is at most
\(q<1\). Shrinking the neighbourhood of \(p_0\) gives
\(\|\mathcal T_p(k_0)-k_0\|<(1-q)r\), where \(r\) is the
ball's radius. Thus \(\mathcal T_p\) maps the closed ball
into itself and is contractive. Its fixed point is the local solution.
Differentiating \(\mathcal R(p,k_*(p))=0\) and solving for
\(\partial_{p_i}k_*\) gives
\eqref{vii:eq:derivada-rama-torsional}.
\end{proof}

The hypothesis is checked on the composed action: Gramian rank
ensures the screen minimum, and the total Hessian controls
the torsional branch. The asserted uniqueness is local around the
specified solution; other branches may coexist in the same
full-rank domain.

\subsection{Effective action and non-affine contribution}
\label{vii:subsec:accion-efectiva-no-afin}

\begin{theorem}[Exact elimination identity]
\label{vii:thm:eliminacion-torsional-no-lineal}
Let \(k_*(p)\) be a stationary branch. Suppose also that the
segment \(t\mapsto (p,tk_*)\), \(0\leq t\leq1\), remains
in the open set where \(F\) is defined. The action correction
relative to the Levi--Civita realization is
\begin{equation}
 \begin{split}
 \Delta S(p)
  &:=Q(p,k_*)+F(p,k_*)-F(p,0)\\
  &=\tfrac14 k_*^{\mathsf T}s(p,k_*)
     -\tfrac12\int_0^1
           k_*^{\mathsf T}s(p,tk_*)\,dt.
 \end{split}
 \label{vii:eq:accion-efectiva-no-lineal}
\end{equation}
Equivalently,
\begin{equation}
 \Delta S=-\tfrac14 k_*^{\mathsf T}s(p,k_*)
  +\tfrac12\int_0^1 k_*^{\mathsf T}
       \bigl[s(p,k_*)-s(p,tk_*)\bigr]\,dt.
 \label{vii:eq:correccion-no-afin-integrada}
\end{equation}
These identities concern the integrated or cellular action.
In a realization with local density, they apply to the density
with its corresponding pairing of forms. The effective action
preserves the declared boundary term separately.
\end{theorem}

\begin{proof}
By the chain rule and
\eqref{vii:eq:covector-torsional-minimo},
\[
 \frac{d}{dt}F(p,tk_*)
      =-\tfrac12 k_*^{\mathsf T}s(p,tk_*).
\]
Integration from \(0\) to \(1\) gives the difference of
material actions. The stationarity equation, evaluated in the
direction \(k_*\), and quadratic homogeneity of \(Q\) give
\[
 2Q(p,k_*)=\tfrac12 k_*^{\mathsf T}s(p,k_*).
\]
Adding both identities proves
\eqref{vii:eq:accion-efectiva-no-lineal}; adding and subtracting
\(\tfrac12k_*^{\mathsf T}s(p,k_*)\) gives
\eqref{vii:eq:correccion-no-afin-integrada}.
\end{proof}

\begin{corollary}[Affine specialization]
\label{vii:cor:recuperacion-eliminacion-afin}
If \(s(p,k)=s(p)\), then
\[
 \Delta S=-\tfrac14 k_*^{\mathsf T}s(p)=-Q(p,k_*).
\]
This is the integrated realization of
\eqref{vii:eq:accion-efectiva-spin}.
\end{corollary}
\begin{proof}
The difference inside the integral in
\eqref{vii:eq:correccion-no-afin-integrada} vanishes. The second
equality follows from the radial stationarity identity.
\end{proof}

\subsection{Metric differential of the eliminated action}
\label{vii:subsec:diferencial-metrico-no-lineal}

\begin{theorem}[Response of the geometric variables]
\label{vii:thm:envolvente-metrica}
On the branch of Theorem~\ref{vii:thm:rama-torsional-local},
\begin{equation}
 \partial_{p_i}\Delta S
 =\tfrac12 k_*^{\mathsf T}(\partial_{p_i}\mathsf A)k_*
  +(\partial_{p_i}F)(p,k_*)-(\partial_{p_i}F)(p,0),
 \label{vii:eq:envolvente-metrica-no-lineal}
\end{equation}
where the partial derivatives hold the coordinates \(k\) fixed.
For fixed Hermitian inner products, each material derivative is calculated by
\begin{equation}
 \begin{split}
 \partial_{p_i}F={}&(\partial_{p_i}\mathfrak a)E
       +\mathfrak a(\partial_{p_i}\chi)\langle b,y\rangle\\
   &+2\mathfrak a\chi\operatorname{Re}
       \langle y,\partial_{p_i}b
                    -(\partial_{p_i}\mathsf C)f_{\min}\rangle.
 \end{split}
 \label{vii:eq:diferencial-geometrico-minimo}
\end{equation}
If \(p_i\) varies the cotetrad, \(\partial_{p_i}\mathsf C\)
includes variation of
\(\mathring\omega(e(p))+K(k)\), the realized routes, and
the trivializations used, according to their declared dependencies.
\end{theorem}

\begin{proof}
The total derivative of \(Q(p,k_*)+F(p,k_*)\) contains, in addition
to its partial derivatives, the term
\[
 \bigl[\partial_kQ(p,k_*)+\partial_kF(p,k_*)\bigr]
                      \partial_{p_i}k_*.
\]
The bracket is zero by the connection equation. Subtracting the derivative
of \(F(p,0)\) proves
\eqref{vii:eq:envolvente-metrica-no-lineal}.
The product rule and
\eqref{vii:eq:diferencial-minimo-completo} prove
\eqref{vii:eq:diferencial-geometrico-minimo}. The composition
\eqref{vii:eq:accion-minima-en-contorsion} requires differentiating
the Levi--Civita connection when \(e\) varies with \(k\) fixed.
\end{proof}

When the positive extension inner product also varies, its
contribution is preserved explicitly. In coordinates, let
\(\mathsf h(p,k)>0\) be its Hermitian matrix and consider the minimum
of \(f^*\mathsf h f\) subject to \(\mathsf C f=b\). Then
\begin{equation}
 \mathsf L=\mathsf C\mathsf h^{-1}\mathsf C^*,\quad
 y=\mathsf L^{-1}b,\quad
 f_{\min}=\mathsf h^{-1}\mathsf C^*y,\quad
 E=\chi b^*y.
 \label{vii:eq:minimo-producto-variable}
\end{equation}
Here \(y\) represents the constraint covector in fixed
bases. The positive matrix \(\mathsf h\) preserves its energetic
role, distinct from that of the Lorentzian metric.

\begin{lemma}[Variation of the extension inner product]
\label{vii:lem:producto-variable-extension}
With the preceding conventions,
\begin{equation}
 \delta E=(\delta\chi)b^*y
   +2\chi\operatorname{Re}\bigl[y^*(\delta b-\delta\mathsf C f_{\min})\bigr]
   +\chi f_{\min}^*(\delta\mathsf h)f_{\min}.
 \label{vii:eq:diferencial-producto-variable}
\end{equation}
\end{lemma}
\begin{proof}
The formula \(\delta\mathsf h^{-1}
=-\mathsf h^{-1}(\delta\mathsf h)\mathsf h^{-1}\) gives
\[
 \delta\mathsf L=(\delta\mathsf C)\mathsf h^{-1}\mathsf C^*
 +\mathsf C\mathsf h^{-1}(\delta\mathsf C)^*
 -\mathsf C\mathsf h^{-1}(\delta\mathsf h)
                              \mathsf h^{-1}\mathsf C^*.
\]
When \(E=\chi b^*\mathsf L^{-1}b\) is differentiated, the term
\(-\chi y^*(\delta\mathsf L)y\) produces the two terms
with \(\delta\mathsf C\) and the positive term with
\(\delta\mathsf h\). This proves the complete identity.
\end{proof}

In a covariant material realization whose variation admits the
metric pairing of
\eqref{vii:eq:fuentes-normalizadas-accion}, the correction determines
\begin{equation}
 \delta_g\Delta S
 =-\tfrac12\int_M\operatorname{vol}_e\,
          \mathcal U^{S,\min}_{\mu\nu}\,\delta g^{\mu\nu},
 \qquad U^{\rm phys,\min}_{\mu\nu}
                      =c\mathcal U^{S,\min}_{\mu\nu}.
 \label{vii:eq:fuente-metrica-minimo-no-lineal}
\end{equation}
Equations \eqref{vii:eq:envolvente-metrica-no-lineal} and
\eqref{vii:eq:diferencial-geometrico-minimo}, supplemented by
\eqref{vii:eq:diferencial-producto-variable} where appropriate,
evaluate its coefficients in the declared metric variations.
The screen minimum, torsional solution, and metric derivative
are three consecutive operations on the same action.
The reduced action preserves any couplings between
cells present in the original routes.

The homogeneous specialization of
\ref{vii:sec:dilucion-rebote} additionally uses the realization of the
tensor obtained: its spatial symmetry, the sign of the contraction
with the observer, and the scaling law. When that evaluation gives
\(\rho_{\rm corr}(a)=-s_0a^{-6}\), the amplitude is read in the
reference state as
\(s_0=-a_{\rm ref}^{6}\rho_{\rm corr}(a_{\rm ref})\).
This reading preserves the material state and its normalization; the
elimination theorem supplies the action to be contracted.

\subsection{Exact non-affine verification model}
\label{vii:subsec:control-no-afin}

The following finite model checks the coefficients and non-affine
dependence of the identities. Its parameters belong to
the algebraic example and keep the problem of
material realization of routes separate.

\begin{proposition}[Exact elimination of a rank-one extension]
\label{vii:prop:ejemplo-no-afin}
Let \(H=\mathbb R^2\), \(Q=\mathbb R\),
\(\mathsf C(k)=(1\ \ k)\), \(b=1\), \(\chi=1\),
\(\mathfrak a=\lambda>0\), and \(Q_{\rm ex}(k)=k^2/2\).
Then
\begin{equation}
 f_{\min}=\frac{(1,k)^{\mathsf T}}{1+k^2},\quad
 F(\lambda,k)=\frac{\lambda}{1+k^2},\quad
 s(\lambda,k)=\frac{4\lambda k}{(1+k^2)^2},\quad
 \mathcal R(\lambda,k)
    =k\left(1-\frac{2\lambda}{(1+k^2)^2}\right).
 \label{vii:eq:ejemplo-minimo-no-afin}
\end{equation}
For \(\lambda=25/2\), the stationary solutions are
\(0,2,-2\). At \(k_*=2\),
\begin{equation}
 \mathsf H_*=\frac{16}{5},\quad s_*=4,\quad
 Q_{\rm ex}(k_*)=2,\quad F(\lambda,k_*)=\frac52,
 \quad \Delta S=-8.
 \label{vii:eq:valores-ejemplo-no-afin}
\end{equation}
The affine expression \(-k_*s_*/4=-2\) differs from the exact
correction. On the positive branch,
\begin{equation}
 \partial_\lambda\Delta S=-\frac45,
 \qquad \partial_\lambda k_* =\frac1{20}
 \quad\text{at }\lambda=25/2.
 \label{vii:eq:derivadas-ejemplo-no-afin}
\end{equation}
\end{proposition}

\begin{proof}
The Gramian is \(1+k^2>0\), so the minimum has the
stated expressions. Differentiating \(F\) gives
\(s=-2\partial_kF\), and then \(\mathcal R=k-s/2\).
For \(\lambda=25/2\), the equation is
\(k[(1+k^2)^2-25]=0\); since \(1+k^2>0\), its roots
are precisely \(0,\pm2\). The total Hessian is
\[
 1-\frac{2\lambda(1-3k^2)}{(1+k^2)^3},
\]
equal to \(16/5\) at both nonzero roots. The correction is
\(2+5/2-25/2=-8\). In the radial identity,
\[
 \int_0^1 k_*s(\lambda,tk_*)\,dt
    =\int_0^1\frac{200t}{(1+4t^2)^2}\,dt=20,
\]
since an antiderivative is \(-25/(1+4t^2)\).
Thus \(\Delta S=8/4-20/2=-8\), including the non-affine term.
For \(\lambda>1/2\), the positive branch satisfies
\(k_*^2=\sqrt{2\lambda}-1\) and
\(\Delta S=\sqrt{2\lambda}-1/2-\lambda\).
Their derivatives give \eqref{vii:eq:derivadas-ejemplo-no-afin}.
The envelope formula gives the same result:
\((\partial_\lambda F)(\lambda,k_*)
 -(\partial_\lambda F)(\lambda,0)=1/5-1=-4/5\).
\end{proof}
